% --- Archon physics formalization source begin ---
% archon:physics
% archon:covers PhyXMiniProblems/problem_phyx_mini_0957.lean
% archon:source-report reports/phyx_mini/problem_phyx_mini_0957.source.json

\chapter{Physics problem phyx\_mini\_0957}
\label{ch:PhyXMiniProblems_problem_phyx_mini_0957}

\paragraph{Problem source.}
\#\# Physical scenario

A particle with charge \textbackslash{}( 2.15\textbackslash{},\textbackslash{}mu\textbackslash{}text\{C\} \textbackslash{}) and mass \textbackslash{}( 3.20 \textbackslash{}times 10\textasciicircum{}\{-11\}\textbackslash{},\textbackslash{}text\{kg\} \textbackslash{}) is initially traveling in the \textbackslash{}( +y \textbackslash{})-direction with a speed \textbackslash{}( v\_0 = 1.45 \textbackslash{}times 10\textasciicircum{}5\textbackslash{},\textbackslash{}text\{m/s\} \textbackslash{}). It then enters a region containing a uniform magnetic field that is directed into, and perpendicular to. The magnitude of the field is \textbackslash{}( 0.420\textbackslash{},\textbackslash{}text\{T\} \textbackslash{}). The region extends a distance of \textbackslash{}( 25.0\textbackslash{},\textbackslash{}text\{cm\} \textbackslash{}) along the initial direction of travel; \textbackslash{}( 75.0\textbackslash{},\textbackslash{}text\{cm\} \textbackslash{}) from the point of entry into the magnetic field region is a wall. The length of the field-free region is thus \textbackslash{}( 50.0\textbackslash{},\textbackslash{}text\{cm\} \textbackslash{}). When the charged particle enters the magnetic field, it follows a curved path whose radius of curvature is \textbackslash{}( R \textbackslash{}). It then leaves the magnetic field after a time \textbackslash{}( t\_1 \textbackslash{}), having been deflected a distance \textbackslash{}( \textbackslash{}Delta x\_1 \textbackslash{}). The particle then travels in the field-free region and strikes the wall after undergoing a total deflection \textbackslash{}( \textbackslash{}Delta x \textbackslash{}).

\#\# Question

Determine \textbackslash{}( \textbackslash{}Delta x \textbackslash{}), the total horizontal deflection.

\#\# Answer choices

- A: A: 0.0245m

- B: B: 0.0305m

- C: C: 0.035m

- D: D: 0.0172m

\#\# Figure caption (auxiliary; use the image as primary evidence)

The image depicts a physics-related diagram involving a charged particle, magnetic field, and two-dimensional motion.

1. **Objects and Elements:**
   - A wall labeled at the top with a gray horizontal bar.
   - A positive charge (represented by a circle with a plus sign) at the bottom.
   - A magnetic field indicated by "X" marks, implying the field is directed into the page.
   - Vectors and trajectories indicated with arrows.

2. **Vectors and Trajectories:**
   - The vector \textbackslash{}( \textbackslash{}mathbf\{v\}\_0 \textbackslash{}) is a green arrow pointing upward from the positive charge, representing the initial velocity.
   - There is a black line segment labeled \textbackslash{}( \textbackslash{}mathbf\{B\} \textbackslash{}) indicating the magnetic field strength.
   - A dashed blue line illustrating the particle's path which first curves and then hits the wall.

3. **Measurements and Labels:**
   - \textbackslash{}( D = 75.0 \textbackslash{}, \textbackslash{}text\{cm\} \textbackslash{}) is the vertical distance between the wall and the initial position of the particle.
   - \textbackslash{}( d = 25.0 \textbackslash{}, \textbackslash{}text\{cm\} \textbackslash{}) is the distance between the launch point and the bottom edge of the field.
   - \textbackslash{}( \textbackslash{}Delta x\_1 \textbackslash{}) is the horizontal distance between the field's edge and the point where the trajectory meets the wall's extension.
   - \textbackslash{}( \textbackslash{}Delta x \textbackslash{}) is the horizontal distance between a perpendicular from point of impact on the wall and a reference line or axis.
   - \textbackslash{}( R \textbackslash{}) is the radius of curvature of the particle’s path within the magnetic field.

4. **Coordinate System:**
   - The x-axis runs horizontally along the bottom, and the y-axis runs vertically along the side. The origin is located at the bottom where the charge starts.

This diagram likely illustrates the motion of a charged particle in a magnetic field and its interaction with an external barrier (the wall).

\#\# Dataset metadata

Magnetism; Physical Model Grounding Reasoning, Spatial Relation Reasoning

\paragraph{Recorded answer/context.}
B: B: 0.0305m

\paragraph{Figure/image path.}
/root/proposal\_for\_physic/science-mango/phyx\_mini\_run/phyx\_data/test\_image/957.png

\paragraph{Formalization target.}
create a compiling Lean file with sorry bodies at `PhyXMiniProblems/problem\_phyx\_mini\_0957.lean`. The Lean declarations must preserve the physical quantities, dimensions or dimensional roles, figure labels, governing-law hypotheses, and final relation expressed by this problem.
Use Mathlib/Physlib names found through LeanExplore where available. If a physics API is missing, introduce faithful local abstractions rather than scalar placeholder aliases.

\begin{theorem}[Physics formalization target]
\label{thm:physics:phyx_mini_0957:target}
\lean{PhyXMiniProblems.ProblemPhyXMini0957.problem_phyx_mini_0957}
\uses{def:physics:phyx-mini-0957:phyxminiproblems-problemphyxmini0957-lengthinmeters, def:physics:phyx-mini-0957:phyxminiproblems-problemphyxmini0957-massinkilograms, def:physics:phyx-mini-0957:phyxminiproblems-problemphyxmini0957-chargemagnitudeincoulombs, def:physics:phyx-mini-0957:phyxminiproblems-problemphyxmini0957-speedinmeterspersecond, def:physics:phyx-mini-0957:phyxminiproblems-problemphyxmini0957-magneticfluxdensityinteslas, def:physics:phyx-mini-0957:phyxminiproblems-problemphyxmini0957-chargedparticledeflectionsetup, def:physics:phyx-mini-0957:phyxminiproblems-problemphyxmini0957-matcheswrittenproblemdata, def:physics:phyx-mini-0957:phyxminiproblems-problemphyxmini0957-matchessuppliedmagneticdeflectionfigure, def:physics:phyx-mini-0957:phyxminiproblems-problemphyxmini0957-hasphysicaldeflectionparameters, def:physics:phyx-mini-0957:phyxminiproblems-problemphyxmini0957-satisfiesfieldandwallgeometry, def:physics:phyx-mini-0957:phyxminiproblems-problemphyxmini0957-hasuniformfinitemagneticfield, def:physics:phyx-mini-0957:phyxminiproblems-problemphyxmini0957-satisfiesperpendicularmagneticcircularmotion, def:physics:phyx-mini-0957:phyxminiproblems-problemphyxmini0957-satisfiesfieldfreetangentmotion, def:physics:phyx-mini-0957:phyxminiproblems-problemphyxmini0957-recordeddatasetanswer, def:physics:phyx-mini-0957:phyxminiproblems-problemphyxmini0957-isuniqueclosestdisplayeddeflection}
The assigned autoformalize agent should translate this physics problem into Lean declarations in the covered file, with theorem and lemma proof bodies written as `by sorry`.
\end{theorem}
\begin{proof}
This is an autoformalization task, not a proof task. Produce statements that can later be proved by the physics prover without weakening the physical model.
\end{proof}

% --- Archon declaration topology begin ---
\section{Declaration topology}
\begin{definition}[magnetic Flux Density Dimension]
  \label{def:physics:phyx-mini-0957:phyxminiproblems-problemphyxmini0957-magneticfluxdensitydimension}
  \lean{PhyXMiniProblems.ProblemPhyXMini0957.magneticFluxDensityDimension}
  Magnetic flux density has coherent-SI dimension M T⁻¹ C⁻¹ (tesla)
\end{definition}
\begin{proof}
  This is the declared model definition of magnetic Flux Density Dimension.
\end{proof}
\begin{definition}[Length Quantity]
  \label{def:physics:phyx-mini-0957:phyxminiproblems-problemphyxmini0957-lengthquantity}
  \lean{PhyXMiniProblems.ProblemPhyXMini0957.LengthQuantity}
  A nonnegative, unit-independent physical length
\end{definition}
\begin{proof}
  This is the declared model definition of Length Quantity.
\end{proof}
\begin{definition}[Time Quantity]
  \label{def:physics:phyx-mini-0957:phyxminiproblems-problemphyxmini0957-timequantity}
  \lean{PhyXMiniProblems.ProblemPhyXMini0957.TimeQuantity}
  A nonnegative, unit-independent duration
\end{definition}
\begin{proof}
  This is the declared model definition of Time Quantity.
\end{proof}
\begin{definition}[Mass Quantity]
  \label{def:physics:phyx-mini-0957:phyxminiproblems-problemphyxmini0957-massquantity}
  \lean{PhyXMiniProblems.ProblemPhyXMini0957.MassQuantity}
  A nonnegative, unit-independent mass
\end{definition}
\begin{proof}
  This is the declared model definition of Mass Quantity.
\end{proof}
\begin{definition}[Charge Magnitude Quantity]
  \label{def:physics:phyx-mini-0957:phyxminiproblems-problemphyxmini0957-chargemagnitudequantity}
  \lean{PhyXMiniProblems.ProblemPhyXMini0957.ChargeMagnitudeQuantity}
  A nonnegative, unit-independent electric-charge magnitude
\end{definition}
\begin{proof}
  This is the declared model definition of Charge Magnitude Quantity.
\end{proof}
\begin{definition}[Speed Quantity]
  \label{def:physics:phyx-mini-0957:phyxminiproblems-problemphyxmini0957-speedquantity}
  \lean{PhyXMiniProblems.ProblemPhyXMini0957.SpeedQuantity}
  A nonnegative, unit-independent speed
\end{definition}
\begin{proof}
  This is the declared model definition of Speed Quantity.
\end{proof}
\begin{definition}[Magnetic Flux Density Quantity]
  \label{def:physics:phyx-mini-0957:phyxminiproblems-problemphyxmini0957-magneticfluxdensityquantity}
  \lean{PhyXMiniProblems.ProblemPhyXMini0957.MagneticFluxDensityQuantity}
  \uses{def:physics:phyx-mini-0957:phyxminiproblems-problemphyxmini0957-magneticfluxdensitydimension}
  A nonnegative, unit-independent magnetic-flux-density magnitude
\end{definition}
\begin{proof}
  This is the declared model definition of Magnetic Flux Density Quantity.
\end{proof}
\begin{definition}[length Readout]
  \label{def:physics:phyx-mini-0957:phyxminiproblems-problemphyxmini0957-lengthreadout}
  \lean{PhyXMiniProblems.ProblemPhyXMini0957.lengthReadout}
  \uses{def:physics:phyx-mini-0957:phyxminiproblems-problemphyxmini0957-lengthquantity}
  Read a physical length in a selected Physlib length unit
\end{definition}
\begin{proof}
  This is the declared model definition of length Readout.
\end{proof}
\begin{definition}[length In Meters]
  \label{def:physics:phyx-mini-0957:phyxminiproblems-problemphyxmini0957-lengthinmeters}
  \lean{PhyXMiniProblems.ProblemPhyXMini0957.lengthInMeters}
  \uses{def:physics:phyx-mini-0957:phyxminiproblems-problemphyxmini0957-lengthquantity, def:physics:phyx-mini-0957:phyxminiproblems-problemphyxmini0957-lengthreadout}
  Read a physical length in metres
\end{definition}
\begin{proof}
  This is the declared model definition of length In Meters.
\end{proof}
\begin{definition}[length In Centimeters]
  \label{def:physics:phyx-mini-0957:phyxminiproblems-problemphyxmini0957-lengthincentimeters}
  \lean{PhyXMiniProblems.ProblemPhyXMini0957.lengthInCentimeters}
  \uses{def:physics:phyx-mini-0957:phyxminiproblems-problemphyxmini0957-lengthquantity, def:physics:phyx-mini-0957:phyxminiproblems-problemphyxmini0957-lengthreadout}
  Read a physical length in centimetres
\end{definition}
\begin{proof}
  This is the declared model definition of length In Centimeters.
\end{proof}
\begin{definition}[time In Seconds]
  \label{def:physics:phyx-mini-0957:phyxminiproblems-problemphyxmini0957-timeinseconds}
  \lean{PhyXMiniProblems.ProblemPhyXMini0957.timeInSeconds}
  \uses{def:physics:phyx-mini-0957:phyxminiproblems-problemphyxmini0957-timequantity}
  Read a duration in coherent-SI seconds
\end{definition}
\begin{proof}
  This is the declared model definition of time In Seconds.
\end{proof}
\begin{definition}[mass In Kilograms]
  \label{def:physics:phyx-mini-0957:phyxminiproblems-problemphyxmini0957-massinkilograms}
  \lean{PhyXMiniProblems.ProblemPhyXMini0957.massInKilograms}
  \uses{def:physics:phyx-mini-0957:phyxminiproblems-problemphyxmini0957-massquantity}
  Read a mass in coherent-SI kilograms
\end{definition}
\begin{proof}
  This is the declared model definition of mass In Kilograms.
\end{proof}
\begin{definition}[charge Magnitude In Coulombs]
  \label{def:physics:phyx-mini-0957:phyxminiproblems-problemphyxmini0957-chargemagnitudeincoulombs}
  \lean{PhyXMiniProblems.ProblemPhyXMini0957.chargeMagnitudeInCoulombs}
  \uses{def:physics:phyx-mini-0957:phyxminiproblems-problemphyxmini0957-chargemagnitudequantity}
  Read an electric-charge magnitude in coherent-SI coulombs
\end{definition}
\begin{proof}
  This is the declared model definition of charge Magnitude In Coulombs.
\end{proof}
\begin{definition}[speed In Meters Per Second]
  \label{def:physics:phyx-mini-0957:phyxminiproblems-problemphyxmini0957-speedinmeterspersecond}
  \lean{PhyXMiniProblems.ProblemPhyXMini0957.speedInMetersPerSecond}
  \uses{def:physics:phyx-mini-0957:phyxminiproblems-problemphyxmini0957-speedquantity}
  Read a speed in coherent-SI metres per second
\end{definition}
\begin{proof}
  This is the declared model definition of speed In Meters Per Second.
\end{proof}
\begin{definition}[magnetic Flux Density In Teslas]
  \label{def:physics:phyx-mini-0957:phyxminiproblems-problemphyxmini0957-magneticfluxdensityinteslas}
  \lean{PhyXMiniProblems.ProblemPhyXMini0957.magneticFluxDensityInTeslas}
  \uses{def:physics:phyx-mini-0957:phyxminiproblems-problemphyxmini0957-magneticfluxdensityquantity}
  Read a magnetic-flux-density magnitude in coherent-SI teslas
\end{definition}
\begin{proof}
  This is the declared model definition of magnetic Flux Density In Teslas.
\end{proof}
\begin{definition}[Charge Sign]
  \label{def:physics:phyx-mini-0957:phyxminiproblems-problemphyxmini0957-chargesign}
  \lean{PhyXMiniProblems.ProblemPhyXMini0957.ChargeSign}
  Sign of the particle's electric charge
\end{definition}
\begin{proof}
  This is the declared model definition of Charge Sign.
\end{proof}
\begin{definition}[Planar Direction]
  \label{def:physics:phyx-mini-0957:phyxminiproblems-problemphyxmini0957-planardirection}
  \lean{PhyXMiniProblems.ProblemPhyXMini0957.PlanarDirection}
  Named directions in the plane of image 957.png
\end{definition}
\begin{proof}
  This is the declared model definition of Planar Direction.
\end{proof}
\begin{definition}[Page Direction]
  \label{def:physics:phyx-mini-0957:phyxminiproblems-problemphyxmini0957-pagedirection}
  \lean{PhyXMiniProblems.ProblemPhyXMini0957.PageDirection}
  Literal arrow directions on the two-dimensional page
\end{definition}
\begin{proof}
  This is the declared model definition of Page Direction.
\end{proof}
\begin{definition}[Page Normal Direction]
  \label{def:physics:phyx-mini-0957:phyxminiproblems-problemphyxmini0957-pagenormaldirection}
  \lean{PhyXMiniProblems.ProblemPhyXMini0957.PageNormalDirection}
  Directions normal to the page, represented by dots or crosses
\end{definition}
\begin{proof}
  This is the declared model definition of Page Normal Direction.
\end{proof}
\begin{definition}[Relative Orientation]
  \label{def:physics:phyx-mini-0957:phyxminiproblems-problemphyxmini0957-relativeorientation}
  \lean{PhyXMiniProblems.ProblemPhyXMini0957.RelativeOrientation}
  Relative directions needed by the perpendicular magnetic-motion model
\end{definition}
\begin{proof}
  This is the declared model definition of Relative Orientation.
\end{proof}
\begin{definition}[Magnetic Field Model]
  \label{def:physics:phyx-mini-0957:phyxminiproblems-problemphyxmini0957-magneticfieldmodel}
  \lean{PhyXMiniProblems.ProblemPhyXMini0957.MagneticFieldModel}
  Idealized spatial models for the applied field
\end{definition}
\begin{proof}
  This is the declared model definition of Magnetic Field Model.
\end{proof}
\begin{definition}[Figure Quantity Label]
  \label{def:physics:phyx-mini-0957:phyxminiproblems-problemphyxmini0957-figurequantitylabel}
  \lean{PhyXMiniProblems.ProblemPhyXMini0957.FigureQuantityLabel}
  Symbolic physical-quantity labels visible in the supplied raster
\end{definition}
\begin{proof}
  This is the declared model definition of Figure Quantity Label.
\end{proof}
\begin{definition}[expected Printed Symbol]
  \label{def:physics:phyx-mini-0957:phyxminiproblems-problemphyxmini0957-expectedprintedsymbol}
  \lean{PhyXMiniProblems.ProblemPhyXMini0957.expectedPrintedSymbol}
  \uses{def:physics:phyx-mini-0957:phyxminiproblems-problemphyxmini0957-figurequantitylabel}
  Expected printed symbol for each physical-quantity label
\end{definition}
\begin{proof}
  This is the declared model definition of expected Printed Symbol.
\end{proof}
\begin{definition}[Magnetic Deflection Figure]
  \label{def:physics:phyx-mini-0957:phyxminiproblems-problemphyxmini0957-magneticdeflectionfigure}
  \lean{PhyXMiniProblems.ProblemPhyXMini0957.MagneticDeflectionFigure}
  \uses{def:physics:phyx-mini-0957:phyxminiproblems-problemphyxmini0957-lengthquantity, def:physics:phyx-mini-0957:phyxminiproblems-problemphyxmini0957-pagedirection, def:physics:phyx-mini-0957:phyxminiproblems-problemphyxmini0957-pagenormaldirection, def:physics:phyx-mini-0957:phyxminiproblems-problemphyxmini0957-figurequantitylabel}
  Literal presentation data from the primary raster 957.png. The two printed numerical length labels remain physical quantities rather than bare scalars
\end{definition}
\begin{proof}
  This is the declared model definition of Magnetic Deflection Figure.
\end{proof}
\begin{definition}[Charged Particle Deflection Setup]
  \label{def:physics:phyx-mini-0957:phyxminiproblems-problemphyxmini0957-chargedparticledeflectionsetup}
  \lean{PhyXMiniProblems.ProblemPhyXMini0957.ChargedParticleDeflectionSetup}
  \uses{def:physics:phyx-mini-0957:phyxminiproblems-problemphyxmini0957-lengthquantity, def:physics:phyx-mini-0957:phyxminiproblems-problemphyxmini0957-timequantity, def:physics:phyx-mini-0957:phyxminiproblems-problemphyxmini0957-massquantity, def:physics:phyx-mini-0957:phyxminiproblems-problemphyxmini0957-chargemagnitudequantity, def:physics:phyx-mini-0957:phyxminiproblems-problemphyxmini0957-speedquantity, def:physics:phyx-mini-0957:phyxminiproblems-problemphyxmini0957-magneticfluxdensityquantity, def:physics:phyx-mini-0957:phyxminiproblems-problemphyxmini0957-chargesign, def:physics:phyx-mini-0957:phyxminiproblems-problemphyxmini0957-planardirection, def:physics:phyx-mini-0957:phyxminiproblems-problemphyxmini0957-pagenormaldirection, def:physics:phyx-mini-0957:phyxminiproblems-problemphyxmini0957-relativeorientation, def:physics:phyx-mini-0957:phyxminiproblems-problemphyxmini0957-magneticfieldmodel, def:physics:phyx-mini-0957:phyxminiproblems-problemphyxmini0957-magneticdeflectionfigure}
  Independent physical quantities and observables in the experiment. In particular, totalHorizontalDeflection is not defined from the target formula or from an answer choice
\end{definition}
\begin{proof}
  This is the declared model definition of Charged Particle Deflection Setup.
\end{proof}
\begin{definition}[Matches Written Problem Data]
  \label{def:physics:phyx-mini-0957:phyxminiproblems-problemphyxmini0957-matcheswrittenproblemdata}
  \lean{PhyXMiniProblems.ProblemPhyXMini0957.MatchesWrittenProblemData}
  \uses{def:physics:phyx-mini-0957:phyxminiproblems-problemphyxmini0957-lengthincentimeters, def:physics:phyx-mini-0957:phyxminiproblems-problemphyxmini0957-massinkilograms, def:physics:phyx-mini-0957:phyxminiproblems-problemphyxmini0957-chargemagnitudeincoulombs, def:physics:phyx-mini-0957:phyxminiproblems-problemphyxmini0957-speedinmeterspersecond, def:physics:phyx-mini-0957:phyxminiproblems-problemphyxmini0957-magneticfluxdensityinteslas, def:physics:phyx-mini-0957:phyxminiproblems-problemphyxmini0957-chargedparticledeflectionsetup}
  Numerical and qualitative data explicitly stated in the problem prose
\end{definition}
\begin{proof}
  This is the declared model definition of Matches Written Problem Data.
\end{proof}
\begin{definition}[Matches Supplied Magnetic Deflection Figure]
  \label{def:physics:phyx-mini-0957:phyxminiproblems-problemphyxmini0957-matchessuppliedmagneticdeflectionfigure}
  \lean{PhyXMiniProblems.ProblemPhyXMini0957.MatchesSuppliedMagneticDeflectionFigure}
  \uses{def:physics:phyx-mini-0957:phyxminiproblems-problemphyxmini0957-expectedprintedsymbol, def:physics:phyx-mini-0957:phyxminiproblems-problemphyxmini0957-chargedparticledeflectionsetup}
  Primary-raster evidence. The image places positive x to the left, positive y upward, shows crosses for an into-page field, and depicts the circular arc joining a straight tangent at the field boundary
\end{definition}
\begin{proof}
  This is the declared model definition of Matches Supplied Magnetic Deflection Figure.
\end{proof}
\begin{definition}[Has Physical Deflection Parameters]
  \label{def:physics:phyx-mini-0957:phyxminiproblems-problemphyxmini0957-hasphysicaldeflectionparameters}
  \lean{PhyXMiniProblems.ProblemPhyXMini0957.HasPhysicalDeflectionParameters}
  \uses{def:physics:phyx-mini-0957:phyxminiproblems-problemphyxmini0957-lengthinmeters, def:physics:phyx-mini-0957:phyxminiproblems-problemphyxmini0957-timeinseconds, def:physics:phyx-mini-0957:phyxminiproblems-problemphyxmini0957-massinkilograms, def:physics:phyx-mini-0957:phyxminiproblems-problemphyxmini0957-chargemagnitudeincoulombs, def:physics:phyx-mini-0957:phyxminiproblems-problemphyxmini0957-speedinmeterspersecond, def:physics:phyx-mini-0957:phyxminiproblems-problemphyxmini0957-magneticfluxdensityinteslas, def:physics:phyx-mini-0957:phyxminiproblems-problemphyxmini0957-chargedparticledeflectionsetup}
  Positivity and nondegeneracy conditions for the depicted orbit
\end{definition}
\begin{proof}
  This is the declared model definition of Has Physical Deflection Parameters.
\end{proof}
\begin{definition}[Satisfies Field And Wall Geometry]
  \label{def:physics:phyx-mini-0957:phyxminiproblems-problemphyxmini0957-satisfiesfieldandwallgeometry}
  \lean{PhyXMiniProblems.ProblemPhyXMini0957.SatisfiesFieldAndWallGeometry}
  \uses{def:physics:phyx-mini-0957:phyxminiproblems-problemphyxmini0957-lengthinmeters, def:physics:phyx-mini-0957:phyxminiproblems-problemphyxmini0957-chargedparticledeflectionsetup}
  The vertical distance from the field exit to the wall is D - d
\end{definition}
\begin{proof}
  This is the declared model definition of Satisfies Field And Wall Geometry.
\end{proof}
\begin{definition}[Has Uniform Finite Magnetic Field]
  \label{def:physics:phyx-mini-0957:phyxminiproblems-problemphyxmini0957-hasuniformfinitemagneticfield}
  \lean{PhyXMiniProblems.ProblemPhyXMini0957.HasUniformFiniteMagneticField}
  \uses{def:physics:phyx-mini-0957:phyxminiproblems-problemphyxmini0957-magneticfluxdensityinteslas, def:physics:phyx-mini-0957:phyxminiproblems-problemphyxmini0957-chargedparticledeflectionsetup}
  The scalar flux-density magnitude calibrates Physlib's magnetic vector field throughout the field strip, while the named field-free region has zero field
\end{definition}
\begin{proof}
  This is the declared model definition of Has Uniform Finite Magnetic Field.
\end{proof}
\begin{definition}[Satisfies Perpendicular Magnetic Circular Motion]
  \label{def:physics:phyx-mini-0957:phyxminiproblems-problemphyxmini0957-satisfiesperpendicularmagneticcircularmotion}
  \lean{PhyXMiniProblems.ProblemPhyXMini0957.SatisfiesPerpendicularMagneticCircularMotion}
  \uses{def:physics:phyx-mini-0957:phyxminiproblems-problemphyxmini0957-lengthinmeters, def:physics:phyx-mini-0957:phyxminiproblems-problemphyxmini0957-timeinseconds, def:physics:phyx-mini-0957:phyxminiproblems-problemphyxmini0957-massinkilograms, def:physics:phyx-mini-0957:phyxminiproblems-problemphyxmini0957-chargemagnitudeincoulombs, def:physics:phyx-mini-0957:phyxminiproblems-problemphyxmini0957-speedinmeterspersecond, def:physics:phyx-mini-0957:phyxminiproblems-problemphyxmini0957-magneticfluxdensityinteslas, def:physics:phyx-mini-0957:phyxminiproblems-problemphyxmini0957-chargedparticledeflectionsetup}
  School-physics laws for the positively charged particle while its velocity is perpendicular to the uniform magnetic field. The angle θ is measured from the initial +y direction toward displayed +x
\end{definition}
\begin{proof}
  This is the declared model definition of Satisfies Perpendicular Magnetic Circular Motion.
\end{proof}
\begin{definition}[Satisfies Field Free Tangent Motion]
  \label{def:physics:phyx-mini-0957:phyxminiproblems-problemphyxmini0957-satisfiesfieldfreetangentmotion}
  \lean{PhyXMiniProblems.ProblemPhyXMini0957.SatisfiesFieldFreeTangentMotion}
  \uses{def:physics:phyx-mini-0957:phyxminiproblems-problemphyxmini0957-lengthinmeters, def:physics:phyx-mini-0957:phyxminiproblems-problemphyxmini0957-chargedparticledeflectionsetup}
  After leaving the field, the particle retains the velocity tangent to the circular arc. Consequently its extra horizontal displacement across the field-free vertical distance is ℓ tan θ. This is a governing motion law, not the requested eliminated formula or a numerical answer
\end{definition}
\begin{proof}
  This is the declared model definition of Satisfies Field Free Tangent Motion.
\end{proof}
\begin{definition}[Answer Choice]
  \label{def:physics:phyx-mini-0957:phyxminiproblems-problemphyxmini0957-answerchoice}
  \lean{PhyXMiniProblems.ProblemPhyXMini0957.AnswerChoice}
  The four answer labels printed with the source problem
\end{definition}
\begin{proof}
  This is the declared model definition of Answer Choice.
\end{proof}
\begin{definition}[displayed Deflection In Meters]
  \label{def:physics:phyx-mini-0957:phyxminiproblems-problemphyxmini0957-answerchoice-displayeddeflectioninmeters}
  \lean{PhyXMiniProblems.ProblemPhyXMini0957.AnswerChoice.displayedDeflectionInMeters}
  \uses{def:physics:phyx-mini-0957:phyxminiproblems-problemphyxmini0957-answerchoice}
  Horizontal deflection printed beside each choice, in metres
\end{definition}
\begin{proof}
  This is the declared model definition of displayed Deflection In Meters.
\end{proof}
\begin{definition}[recorded Dataset Answer]
  \label{def:physics:phyx-mini-0957:phyxminiproblems-problemphyxmini0957-recordeddatasetanswer}
  \lean{PhyXMiniProblems.ProblemPhyXMini0957.recordedDatasetAnswer}
  \uses{def:physics:phyx-mini-0957:phyxminiproblems-problemphyxmini0957-answerchoice}
  The answer label recorded by the source dataset
\end{definition}
\begin{proof}
  This is the declared model definition of recorded Dataset Answer.
\end{proof}
\begin{definition}[displayed Deflection Error]
  \label{def:physics:phyx-mini-0957:phyxminiproblems-problemphyxmini0957-displayeddeflectionerror}
  \lean{PhyXMiniProblems.ProblemPhyXMini0957.displayedDeflectionError}
  \uses{def:physics:phyx-mini-0957:phyxminiproblems-problemphyxmini0957-lengthinmeters, def:physics:phyx-mini-0957:phyxminiproblems-problemphyxmini0957-chargedparticledeflectionsetup, def:physics:phyx-mini-0957:phyxminiproblems-problemphyxmini0957-answerchoice}
  Absolute discrepancy between a physical deflection and a displayed choice
\end{definition}
\begin{proof}
  This is the declared model definition of displayed Deflection Error.
\end{proof}
\begin{definition}[Is Unique Closest Displayed Deflection]
  \label{def:physics:phyx-mini-0957:phyxminiproblems-problemphyxmini0957-isuniqueclosestdisplayeddeflection}
  \lean{PhyXMiniProblems.ProblemPhyXMini0957.IsUniqueClosestDisplayedDeflection}
  \uses{def:physics:phyx-mini-0957:phyxminiproblems-problemphyxmini0957-chargedparticledeflectionsetup, def:physics:phyx-mini-0957:phyxminiproblems-problemphyxmini0957-answerchoice, def:physics:phyx-mini-0957:phyxminiproblems-problemphyxmini0957-displayeddeflectionerror}
  A choice is the unique nearest displayed value to the physical deflection. This comparison is appropriate for the source's multiple-choice decimal approximations and does not assert that a rounded display is an exact length
\end{definition}
\begin{proof}
  This is the declared model definition of Is Unique Closest Displayed Deflection.
\end{proof}
\begin{lemma}[curvature Radius eq mass mul speed div charge mul field]
  \label{lem:physics:phyx-mini-0957:phyxminiproblems-problemphyxmini0957-curvatureradius-eq-mass-mul-speed-div-charge-mul-field}
  \lean{PhyXMiniProblems.ProblemPhyXMini0957.curvatureRadius_eq_mass_mul_speed_div_charge_mul_field}
  \uses{def:physics:phyx-mini-0957:phyxminiproblems-problemphyxmini0957-lengthinmeters, def:physics:phyx-mini-0957:phyxminiproblems-problemphyxmini0957-massinkilograms, def:physics:phyx-mini-0957:phyxminiproblems-problemphyxmini0957-chargemagnitudeincoulombs, def:physics:phyx-mini-0957:phyxminiproblems-problemphyxmini0957-speedinmeterspersecond, def:physics:phyx-mini-0957:phyxminiproblems-problemphyxmini0957-magneticfluxdensityinteslas, def:physics:phyx-mini-0957:phyxminiproblems-problemphyxmini0957-chargedparticledeflectionsetup, def:physics:phyx-mini-0957:phyxminiproblems-problemphyxmini0957-hasphysicaldeflectionparameters, def:physics:phyx-mini-0957:phyxminiproblems-problemphyxmini0957-satisfiesperpendicularmagneticcircularmotion}
  The perpendicular magnetic-orbit law gives R = m v₀ / (q B). This is an intermediate derived relation, not a premise field
\end{lemma}
\begin{proof}
  The conclusion follows from the governing relations and figure readouts named in the statement of curvature Radius eq mass mul speed div charge mul field.
\end{proof}
\begin{lemma}[total Horizontal Deflection eq circle tangent formula]
  \label{lem:physics:phyx-mini-0957:phyxminiproblems-problemphyxmini0957-totalhorizontaldeflection-eq-circle-tangent-formula}
  \lean{PhyXMiniProblems.ProblemPhyXMini0957.totalHorizontalDeflection_eq_circle_tangent_formula}
  \uses{def:physics:phyx-mini-0957:phyxminiproblems-problemphyxmini0957-lengthinmeters, def:physics:phyx-mini-0957:phyxminiproblems-problemphyxmini0957-chargedparticledeflectionsetup, def:physics:phyx-mini-0957:phyxminiproblems-problemphyxmini0957-hasphysicaldeflectionparameters, def:physics:phyx-mini-0957:phyxminiproblems-problemphyxmini0957-satisfiesfieldandwallgeometry, def:physics:phyx-mini-0957:phyxminiproblems-problemphyxmini0957-satisfiesperpendicularmagneticcircularmotion, def:physics:phyx-mini-0957:phyxminiproblems-problemphyxmini0957-satisfiesfieldfreetangentmotion}
  Eliminating the exit angle and the intermediate deflection from the circle and tangent laws gives the total horizontal displacement. The square root is the positive horizontal radius component at the field exit
\end{lemma}
\begin{proof}
  The conclusion follows from the governing relations and figure readouts named in the statement of total Horizontal Deflection eq circle tangent formula.
\end{proof}
\begin{lemma}[total Horizontal Deflection uniquely matches answer B]
  \label{lem:physics:phyx-mini-0957:phyxminiproblems-problemphyxmini0957-totalhorizontaldeflection-uniquely-matches-answerb}
  \lean{PhyXMiniProblems.ProblemPhyXMini0957.totalHorizontalDeflection_uniquely_matches_answerB}
  \uses{def:physics:phyx-mini-0957:phyxminiproblems-problemphyxmini0957-chargedparticledeflectionsetup, def:physics:phyx-mini-0957:phyxminiproblems-problemphyxmini0957-matcheswrittenproblemdata, def:physics:phyx-mini-0957:phyxminiproblems-problemphyxmini0957-hasphysicaldeflectionparameters, def:physics:phyx-mini-0957:phyxminiproblems-problemphyxmini0957-satisfiesfieldandwallgeometry, def:physics:phyx-mini-0957:phyxminiproblems-problemphyxmini0957-satisfiesperpendicularmagneticcircularmotion, def:physics:phyx-mini-0957:phyxminiproblems-problemphyxmini0957-satisfiesfieldfreetangentmotion, def:physics:phyx-mini-0957:phyxminiproblems-problemphyxmini0957-isuniqueclosestdisplayeddeflection}
  The stated numerical data make the physical result approximately 0.03044 m, whose unique nearest displayed value is choice B, 0.0305 m
\end{lemma}
\begin{proof}
  The conclusion follows from the governing relations and figure readouts named in the statement of total Horizontal Deflection uniquely matches answer B.
\end{proof}
% --- Archon declaration topology end ---
% --- Archon physics formalization source end ---
